% ============================================================
%  R. Nielsen, Religionsphilosophie (Kjøbenhavn: Gyldendal, 1869)
%  INDLEDNING  (§§ 1–3, pp. 1–30)
%  TRANSLATION (English)
% ============================================================
\documentclass[12pt,a4paper]{article}
\usepackage[utf8]{inputenc}
\usepackage[T1]{fontenc}
\usepackage[english]{babel}
\usepackage{microtype}
\usepackage{setspace}
\usepackage[margin=1.2in]{geometry}
\usepackage{fancyhdr}
\usepackage{hyperref}

\hypersetup{
  pdftitle={Philosophy of Religion --- Introduction},
  pdfauthor={R. Nielsen},
  hidelinks
}

\pagestyle{fancy}
\fancyhf{}
\rhead{Nielsen, \emph{Philosophy of Religion} (1869) --- Introduction}
\cfoot{\thepage}

\setstretch{1.3}

\title{\textbf{Philosophy of Religion}\\[6pt]
  \large Introduction}
\author{R.\ Nielsen\\[4pt]
  \small Translated from the Danish}
\date{Copenhagen: Gyldendal, 1869, pp.\ 1--30}

\begin{document}
\maketitle
\thispagestyle{fancy}

\bigskip

\begin{center}
\itshape
Faith and knowledge are as principles absolutely heterogeneous.

\smallskip

The heterogeneity of the principles makes no breach in the
universally valid laws of thought, but shows the impossibility
of any scientific transition whatever, either from knowledge to
faith or from faith to knowledge.

\smallskip

The long conflict between religion and science has its original
ground in an obscure conflation of the heterogeneous principles,
and must therefore cease when the conflation ceases.

\smallskip

The proposition that truth in religion is not truth in science
is no longer a paradox, once the absolute boundary of knowledge
coincides with the mystery from which the cognition of faith
proceeds and to which it returns.

\smallskip

It is these principal theses that the \emph{Philosophy of
Religion} in its present form will develop and substantiate.
Whether the development and the substantiation have any
validity, examination must decide.

\bigskip
Copenhagen, 28 December 1868.

R.\ Nielsen.
\end{center}

\bigskip\bigskip

\section*{INTRODUCTION}

The relation of the philosophy of religion to religion cannot be
stated satisfactorily in ordinary conceptual determinations; its
peculiar nature demands a peculiar treatment. What is said
religiously---``Show me your faith by your works''---must be
said philosophically: Show me your principle by its execution.

\subsection*{Religion and Philosophy}
\subsubsection*{§ 1}

If philosophy as a whole is a science, then the philosophy of
religion, insofar as it is to deserve the name of philosophy,
must naturally also be a science; but precisely the distinctive
feature that it is a science \emph{of} religion gives it a
character diverging from the other sciences and stamps it with a
mark of its own. What this means will best be understood by a
glance at the principal forms the philosophy of religion has
assumed in recent times.

The scientific form must be determined by the principle. The
fundamental question is: do religion and philosophy share the
same principle, or are the principles different? If religion's
principle differs from philosophy's, which principle should
science then follow---philosophy's or religion's? Only by
answering these questions can the concept and method of the
philosophy of religion be determined.

If one proceeds from the presupposition that religion and
philosophy share the same principle, the philosophy of religion
becomes a \emph{speculative science}; if one assumes that the
principles are indeed different but that truth must nonetheless
be judged by the principle of knowledge, the philosophy of
religion becomes a \emph{critical science}. If, finally, one
takes the principles to be absolutely heterogeneous and lays the
religious principle as the foundation for the scientific
illumination, the philosophy of religion becomes an
\emph{inverted science}. The first of these standpoints is
represented by Hegel; the second principally by L.\ Feuerbach;
the third is the one brought to application in the present
treatment.

\medskip

\noindent\textit{a) The philosophy of religion as speculative science.}

\medskip

The object both of religion and of philosophy is, according to
Hegel, the eternal truth in its objectivity---God, and nothing
but God, and the explication of God. Philosophy only develops
itself in developing religion, and in developing religion it
develops itself. Thus religion and philosophy coincide into one.
Philosophy is in reality worship, is religion, for it is the
renunciation of subjective notions and opinions in the
occupation with God. Nonetheless the occupation with God in
philosophy is a different one, proceeding in a different way and
manner than in religion. In religion, God is the immediate
object of consciousness; in philosophy, consciousness of God is
God's own consciousness of himself: the difference lies in the
form. Religion as religion is popular, is for all, and for that
very reason is referred to the figurative form of
representation; but such a form is inadequate to the content.
The form is finite, the content infinite: this is a
disproportion. Can this disproportion be overcome? Can the same
content, without altering itself, be presented in two entirely
different forms?

Here lies the difficulty, and on this difficulty Hegel founders.
For wherever truth is at stake as truth, the Idea is so absorbed
into the form that the form cannot possibly be changed without
the content being changed. In the course of its movement of
thought, speculative philosophy of religion is imperceptibly led
to a result different from the one it originally aimed at. The
movement is this: it begins with science needing to acknowledge
and substantiate the truth of religion; the truth of religion is
determined by the true element in the content---that is, by the
greater or lesser quantum of vital ideality that a given
religion, relative to the stage of formal development at which
it finds itself, is capable of expressing; since religion's form
is not scientific, but the scientific form is the only perfect
one, it is brought to consciousness that even the perfect
religion has an imperfect form, and on account of the
imperfection of its form can only express the true in an
imperfect manner. What distinguishes the perfect religion from
philosophy is therefore a form that must be sublated; but in
sublating the form, what is distinctive to religion is sublated,
and therewith religion itself. In losing its form, religion
simultaneously loses its content, for the truth of the content
depends on the form. The upshot, then, is that philosophy, in
order to comprehend what is true in religion, must comprehend
religion itself as an untruth. The negative, repellent relation
between philosophy and religion, seen from science's side, has
found vigorous expression in Neo-Hegelianism.

\medskip

\noindent\textit{b) The philosophy of religion as critical science.}

\medskip

While the speculatively positive tendency holds principally to
the content common to religion and philosophy and strives to
develop this true content of truth into its true form, the
negatively critical tendency---recognizing the principled
difference between religion and philosophy---proceeds principally
from the form peculiar to religion, demonstrates its manifest
conflict with reason, and so arrives through a demonstration of
the untrue in the form at the untrue in the content.

Immediately it looks as though religion had to do with nothing
but actual, if supernatural, facts; but reason developed to
knowledge comprehends that the supernatural facts are not actual,
and the actual ones not supernatural. It is not spirit as spirit,
but the ideal natural forces of the psychic life, that originally
stir in the originally religious consciousness. Religious
representations are indeed not products of conscious
arbitrariness; on the contrary, they are owed to an unconscious
creation, a drive of souls toward self-objectification, a
psychological necessity by which the imagination, overwhelmed by
the Idea, has placed the divine in the human being before the
human being by placing it outside the human being. Religion is
human; it is the first elementary breakthrough of human reason;
it has its origin not in God's love but in human passions.

The dreams of the heart, the movements of the mind, the
longings of the soul contain the basic elements from which the
Idea-inspired imagination has fashioned all the gods of faith,
including the one Almighty. This is confirmed by insight into
the essence both of monotheism and of polytheism. The mighty,
wrathful, vengeful Jehovah of the Old Testament is, seen
scientifically, an idealized reflection of the abstract,
powerful spirit of the Jewish people, of Israel's strong,
law-bound yet self-willed ego. Christ in the New Testament is
not a real God-man, for a real God-man is a contradiction, but
is nonetheless essentially a man-god, the reconciling God of the
wounded heart and the troubled soul, a God for the suffering.
For imagination and feeling the ideal presents itself
immediately, and on that rests not only the wish itself but also
its fulfillment. The human being, with his constraining natural
limitation, wishes for power; if he had power as he has will, he
would feel blessed and free; but he is powerless, and in his
powerlessness he feels the need for divine assistance. Divine
assistance is immediately determined as a God who has power over
nature arbitrarily intervening in the course of things and
helping by a miracle. In the miracle is intuited the
unconstrained will, freed from natural compulsion, a will
capable of annihilating all resistance, satisfying all wishes,
and hearing all prayers. The miracle is the ground of faith, the
way out of hope, the consolation of imagination and the pious
temperament, the holy fairy-tale; when a God has grown so old
and weak that he can no longer perform miracles, he does best to
resign, to cease being God, and likewise a religion that cannot
continually produce miracles must cease to be religion. Here is
the turning-point where the conflict between religion and
philosophy culminates. Religion holds to the miracle, philosophy
to reason. The miracle is in contradiction with reason; when
reason awakes to the cognition of the laws of actuality, the
miracle disappears, and with the miracle religion too.

Thus philosophy, as negative critique, has according to its own
assurance completed the dissolution of religion.

By the negative path, however, a positive result has been
reached. The result is that religion, repudiated and scorned by
``science'', is referred back to its own principle and thereby
invited to rely on its own resources and assert its independence
vis-\`{a}-vis science. The decisive rupture between religion and
science can on these terms benefit religion just as much as
science. The absolute heterogeneity of religion from all
objective knowledge and science has been demonstrated with
penetrating acuity and triumphant superiority by S.\
Kierkegaard.

The fundamental question from which S.\ Kierkegaard sets out is
this: to what extent can truth be taught? Here one encounters
what Socrates in the \emph{Meno} calls a contentious
proposition: that a human being cannot possibly seek what he
knows, and just as little can seek what he does not know---for
what he knows, he cannot seek, since he knows it; and what he
does not know, he cannot seek, since he does not even know what
it is he is to seek. Socrates resolves the difficulty by holding
that all learning and seeking is only a recollecting, so that
the ignorant person needs merely to be reminded in order,
through himself, to call to mind what he knows. Truth is thus
not brought into him but was already in him. This is what we
might call the humanist explanation. From this standpoint
``the moment''---the entry of the divine into time, the
miracle---receives no significance. ``For the Socratic view,
every human being is himself the center, and the whole world
centers only upon him, because his self-knowledge is a knowledge
of God.'' The one who learns the truth learns it through
himself; the teacher is only the occasion; the moment at which
the truth becomes evident is, relative to the truth that is and
that inheres in the essence of the human being, a vanishing and
to that extent an indifferent point in time. ``The point of
departure is a nothing; for in the very moment I discover that I
have known the truth from eternity without knowing it, in that
same instant that moment is hidden in the eternal, absorbed into
it in such a way that I cannot even find it, so to speak, even
if I were to seek it, because here there is no here and there,
but only an \textit{ubique et nusquam}.''

From the exclusively religious standpoint it is otherwise. Here
the human being is presupposed to be, not in truth, but in
untruth. If the one who has lost the truth through the guilt of
the race and through his own guilt is to arrive at the
cognition of truth---namely, that cognition of truth which leads
to blessedness---then it becomes necessary that God not only
enlighten the one in darkness about the untruth, but also, by
himself entering into time and existing as a human being, give
the human being who exists in untruth the condition for
attaining truth. Through God's self-communication in time, the
moment---the determinate point in time---and therewith the
miracle, receives absolutely decisive significance. The miracle
is the paradox of the understanding. By holding fast to the
miracle and positing the paradox, S.\ Kierkegaard has secured
religion against the encroachments of science and erected a
bulwark against the incursion of negative critique; but he has
then also, with decisive aim at his ``introductory problem, not
to Christianity but to becoming a Christian,'' come to rest at
the determination that Christianity is not a doctrine but an
existential communication---an existential communication
expressed in the problem: ``that the individual's eternal
blessedness is decided in time through the relation to something
historical, which is further of such a kind that its composition
includes something that, according to its essence, cannot become
historical and thus must become so by virtue of the absurd.''

According to the negative critique it is science that declares
religion an absurdity and thereby means to dissolve it;
according to Kierkegaard's settlement it is religion that,
despite objective knowledge, declares itself an absurdity and
thereby once and for all excludes science. From this it is seen
that religion and philosophy are not merely heterogeneous but
absolutely heterogeneous. With this presupposition, is a
philosophy of religion possible?

\medskip

\noindent\textit{c) The philosophy of religion as inverted science.}

\medskip

The exclusive relation to God, determined by the miracle, which
is religion's own inescapable presupposition, science cannot
appropriate. This presupposition is the incomprehensible, the
unfathomable, and can only be acknowledged on the condition that
human knowledge acknowledges a barrier, an absolute boundary, at
which it must stop. What in other domains of cognition shows
itself to be contradictory and unreasonable is always something
lower, over which science can freely raise itself; it can
examine this lower, dissolve it, and by its dissolution
continually advance further. But religion---what is
contradictory for knowledge---is not something lower but
something higher, for religion is the highest ideal power of
life; where this power, in virtue of its original essence,
meets one with ``an absurdity,'' the way is objectively
impassable, reason finds the passage blocked, and science must
turn aside.

But what is absurd for scientific thinking---``the paradox''---
nonetheless falls within the thinking consciousness, insofar as
it falls within religion; from which it again follows that what
is unthinkable for science can be thinkable in religion. There
is an essential difference between: overstepping the boundary of
knowledge and overstepping the boundary of thought; between:
falling outside knowledge and falling outside thought. The
distinction between the thinkable and the unthinkable is here to
be determined by the principle from which thinking proceeds.
That the paradox falls outside knowledge is certain; but does it
therefore fall outside thought? For thought to be able to
encounter the paradox, thinking must itself bring the paradox
forward; the paradox is only unthinkable insofar as it is
thought; its dialectical knot is an entanglement of absolutely
heterogeneous threads of thought. Existentially the knot is
insoluble; but it by no means follows from this that it must
also be so intellectually. The intellectual resolution is
conditioned on the heterogeneous threads being unraveled and
separated from one another. Science would not be able to
apprehend and determine its own principle of knowledge if it
were not able, by means of the opposition between what can be
known and what cannot be known, at least negatively to
characterize the principle of faith that is impenetrable to
knowledge. But if it is possible for science to separate the
absolutely heterogeneous principles, then it is also possible
for it to separate the thoughts proceeding from these
principles---possible, that is, to distinguish between thoughts
of knowledge and thoughts of faith, between concepts of
knowledge and concepts of faith.

With the peculiar task of clarifying the relation of tension
between the heterogeneous thoughts, the philosophy of religion
is most precisely to be determined as a \emph{boundary-science}.

The philosophy of religion does not, however, remain standing at
general boundary-determinations for faith and knowledge; it
cannot possibly let religion develop its content without letting
an entire system of concepts of faith come to development. But
the religious content of faith is, according to the principle,
a supra-scientific content; here then enters the contradiction
that a supra-scientific content is taken up into science. If
scientific form were now truly conflated with the content's true
form and introduced and maintained in its place, the
contradiction would certainly be insoluble. But this is by no
means the case. Scientific reflection has precisely the
flexibility that it can bring the form of knowledge to reflect
another and opposite form. For just as bright daylight can bring
the eye to forget the glare and see the objects, so too
reflection can bring thought to see, in knowledge, toward faith
and to forget knowledge. The philosophy of religion is so far
from imposing a foreign form on the religious content that its
formal development consists on the contrary in liberating the
content from the additions of immediately individual
representations, images, feelings, expressions of will, and so
forth, with which it is always intertwined in the life of faith
itself. Insofar as reflection---with regard both to form and to
content---sublates knowledge through knowledge, by aiming to
clarify not a scientific but a supra-scientific truth, the
relation between truth and science is essentially inverted, and
the philosophy of religion becomes in this respect to be
characterized as an \emph{inverted science}.

What is decisive for every science is the grounding. A thorough
development of content can of course be regarded as a peculiar
development of the matter itself---that is, as a development in
which the matter unfolds its moments and grounds itself. But in
all direct sciences, the principle of the matter's grounding is
at bottom one with the general principle of knowledge, so that
the matter's grounding coincides point for point with that
science's own self-grounding. In the inverted science, by
contrast, the grounding is itself inverted. The ground-principle
here is so far from being absorbed into the principle of
knowledge that the two absolutely heterogeneous principles---
those of the miracle and of reason---everywhere repel one
another and in their negative relation determine each other
negatively. About this negative relation, with the
reciprocal determinations corresponding to it, a clear knowledge
can certainly be given; but the general knowledge hovering above
the opposing principles is at the present stage no longer able
to take sides with its own principle. The inverted science has
precisely the task of empowering faith and bringing knowledge
itself to bow before it. As it goes with the grounding
principle, so it must go with the grounding. The inverted
science is to neither suppress nor deny the universally valid
grounds of knowledge, but it is to bring to clear insight how
religion's self-confirmation takes place precisely by the
grounds of knowledge consistently dissolving themselves and
giving way to the absolutely self-validating grounds of faith.
By grounding the conversion and dissolution of the grounds of
knowledge into grounds of faith, the philosophy of religion
grounds its own scientific dissolution, and receives on these
terms significance as a science that sacrifices itself to
religion and sublates itself into doctrine of faith.

\subsection*{Religion and Mythology}
\subsubsection*{§ 2}

The scientific understanding of the opposition between true and
false religion---between religion and mythology---depends on the
relation in which philosophy places itself to religion.
Speculative reason turns mythology into religion; the critical
understanding turns religion into mythology; only in inverted
knowledge can the principled distinction between religion and
mythology be carried through. We shall consider these
standpoints more closely.

\medskip

\noindent\textit{a) Mythology is essentially religion: the speculative view.}

\medskip

With the presupposition that the truth in religion and the truth
in science are not two mutually heterogeneous truths but one and
the same absolute truth, speculative philosophy of religion must
consistently acknowledge religion in every mythology, and regard
the mythical clothing as a husk that can easily be removed.
Religion is spirit, and for spirit as spirit the finite
representations are indifferent. To fathom the essence of spirit
is to fathom the essence of religion; in order to fathom the
essence of religion, philosophy develops religion's concept.
Religion's concept is in its generality a genus-concept, but
according to its particular modes of appearance a
species-concept. The Idea of religion has differentiated its
content and laid out its moments in a plurality of more and less
one-sided religions, in order through these to gather itself out
of dispersal and arrive at completion in the perfect, one, true
religion. ``It has,'' says Hegel, ``been the work of spirit
through millennia to execute the concept of religion and to make
it an object of consciousness.'' The task is to show how the
spirit of humanity, beginning in full naturalness, has been able
step by step to raise itself above the nature-bound condition to
ever greater freedom and finally to reach the turning-point at
which spirit comprehended itself as spirit and expressed the
concept in the exemplary existence that is, in the form of
society, to be appropriated and actualized by all.

Hegel's philosophical-religious systematics is in this respect
characteristic. From nature-religion---the stage at which
consciousness, in a fermenting unity of nature and spirit,
intuits the divine---through the religions of spiritual
individuality or free subjectivity---those of sublimity,
beauty, and purposiveness---the development advances to the
apprehension and determination of the religion that is clear
within itself, the revealed religion corresponding to the
concept of revelation: the absolute religion, Christianity.

According to this view all religions are related, all sprung from
one root, all branches on one trunk. The old distinction between
the nature-religions and revealed religion is in Hegel
overlooked, even to the degree that the ``religion of sublimity''
in Judaism is, for the sake of the concept, placed lower than the
``religion of beauty'' among the Greeks and the ``religion of
purposiveness'' among the Romans. What follows from this? That
the Jewish religion is mythology just as much as the Greek and
Roman! But if Judaism is mythology, then Christianity is also
mythology: the absolute religion has the same origin as the
relative ones. With the absolute religion, then, the mythical
has not disappeared; on the contrary, it has become transparent
to spirit as belonging to religion's distinctive form. Only when
the concept represented in religion is comprehended in knowledge
---that is, when religion is transformed into philosophy---does
the mythical disappear.

``The absolute religion,'' says Hegel, ``is the revealed. Religion
is the revealed only when religion's concept is for itself---that
is, when its concept has become objective to itself, not in a
restricted, finite objectivity, but such that it is according to
its concept objective to itself. We thus have two: consciousness
and the object; but in the religion filled with itself, the
revealed religion that has grasped itself, the content is itself
the object, and this object, the self-knowing being, is
spirit.''\footnote{Hegel, \textit{Religionsphilosophie} II.
Berlin 1840, pp.\ 192--93.}

When religion's content in this manner coincides with
philosophy's Idea, when religious cognition is made one with
speculation's absolute knowledge, who fails to see, even at
first glance, the consequence? Religion demands faith in its
supernatural facts, holds fast to existential immediacy, and
will not let itself be transformed into philosophy; philosophy,
which has transformed the sacred facts into mythical
representations, is confident in knowledge, declares all
positive religion to be mythology, and gives up faith.

\medskip

\noindent\textit{b) Religion is mythology: the critical view.}

\medskip

The conflict between religion and philosophy began among the
Greeks very early, and as for the Romans, two augurs in Cicero's
time could not look at one another without laughing. To be sure,
religion furnished the original basis for spiritual life among
the Greeks as well as the Romans; but in both it is also shown
that the form of religious representations satisfies spirit only
at the first, childlike stage of popular life: the developmental
law is that consciousness, with increasing reflection and
advancing culture, outgrows religion. When Socrates had
discovered the dialectic of subjectivity, set self-knowledge as
the goal of philosophy, and to that extent introduced ``new
gods,'' popular religion could no longer withstand the attacks
of the philosophers. Even if the concept of myth was not in
antiquity developed with the clarity and depth achieved in
modern times, it nonetheless became, as regards the lives of the
gods and the deeds of the heroes, soon sufficiently evident
that---to use an expression of Strauss---``the divine could not
have happened in such a way that what thus happened could not
be divine.'' This dissolution, which explains the myths but
annihilates the gods, is not a spiritual regression but, in the
name of thought, reason, and self-consciousness, an essential
progress. The mythological is the unconsciously illusory; but
this openly illusory element religion cannot shed without
exposing its innate weakness and surrendering itself. The
supernatural phenomena in the New Testament are just as contrary
to reason as the events in the Greek war of the gods. Religion
and mythology are synonymous concepts, and the difference
between lower and higher, more imperfect and more perfect
religions only a difference between lower and higher, more
imperfect and more perfect mythologies.

The higher a mythology stands, the more filled it is with
substance and spirit, the longer it will naturally hold out
against the consciousness of culture; but precisely the
circumstance that the most spirit-filled of all myths---the myth
of the Risen One---has been able to carry culture through
eighteen centuries and hold out against the spirit of humanity,
until consciousness at last got the better of all obscure
presuppositions and came to itself, becoming ``transparent to
itself as the free, self-subsistent universal''---this shows
what it is that in the ``revealed'' religion becomes the
revealed. It is not a fundamental difference between myth and
revelation that comes to light; no, what comes to light now is
something entirely different: it is the great objective truth
that the spirit of humanity---not the spirit of the individual
human being, but the spirit of humanity in its principle---is
its own divinity, is in the religious self-doubling both subject
and object, as finite spirit subject, as spirit in inner
infinity object. This is no blasphemous self-deification; it is
the discovery of the Absolute, of the Idea as Idea, of the
divine-human unity of the subjective and the objective.

By carrying through the principle of knowledge, philosophical
critique has consistently dissolved religion as religion---for
to want to preserve what is religious when religion's principle
has been dissolved is of course nonsense---and has thereby at
the same time reached the boundary at which the powers of life
react: what is religious asserts its right, faith stands up
against knowledge, and the heterogeneous principles measure
themselves against one another.

\medskip

\noindent\textit{c) Religion is absolutely heterogeneous from mythology: the anti-rationalist view.}

\medskip

There is in both the critical and the speculative philosophy's
stance toward religion an ambiguity that must from first to last
blind the observer and confuse his apprehension. This ambiguity
lies in the use of the word \emph{truth}. That the eternal,
absolute truth is only one, and that the one divine truth cannot
have become at odds with itself or have divided itself into two
mutually conflicting truths, is certainly as unambiguous as it
is incontestable a presupposition. But when one then simply
overlooks the difference between a divine and a human knowledge
of truth---overlooks that what is absolute in God's knowledge
must be relative in the human being's, overlooks the principled
difference between religious and scientific cognition---the
ambiguity arises. It looks as though philosophy were truly
entering into religion's own presuppositions, letting these
presuppositions develop their religious consequences, and thereby
arriving at conflicting, mutually canceling results. Yes, if
things truly proceeded in that way, it could truthfully be said
that philosophical critique let religion itself complete its own
dissolution. But things do not proceed in that way. Instead of
beginning with the religious principle, science begins on the
contrary with the principle of knowledge. When applied to
religion, the principle of knowledge becomes consistently a
religion-negating principle. When it then looks as though
critique proves religion's self-dissolution, this is only an
appearance. The truth is that critique gets out what it has
itself put in; it begins by putting in the negation of religion
and ends by getting the negation of religion out: the dissolving
critique moves in a circle. For what is it that science sets up
in principle as the mythical? It is precisely the relation to
God distinctive of religion. To clarify and determine the
relation to God is to clarify and determine the fundamental
conditions decisive for spiritual life. Now science ought to
know, however, that these fundamental conditions present
themselves differently in the domain of faith than in the
domain of knowledge. If it is truly science's serious intent to
get spiritual life itself clarified and understood from the
ground up, then it must, at least for the time being, be capable
of showing the self-denial of seeing past its own principle and
letting religion have a hearing.

That the principle of religion is absolutely heterogeneous from
the principle of science can be seen from the grounding. All
grounds of knowledge are in the deepest sense grounds of
necessity; the grounds of religion, by contrast, are grounds of
will. God's will is the Absolute of religion. If it were a
contradiction---as the thoughtless assert---that science should
come to know a principle absolutely heterogeneous from knowledge,
then it would have to be precisely the negative critique that has
deluded itself into the contradiction. What does this critique
want? To dissolve positive religion by means of knowledge. Why
does it want to dissolve positive religion? Because its original
presuppositions are absolutely heterogeneous from science.

\end{document}
